We then consider the canonical morphism (2.7.2):
@@M0@@
let @@M1@@ be the cone of this morphism. It is again a @@M2@@-perfect complex: indeed, for every @@M3@@, @@M4@@ is a perfect complex on @@M5@@, since this holds when @@M6@@ is bounded with terms locally free of finite type over @@M7@@; consequently, the complex @@M8@@ is @@M9@@-perfect, and the same therefore holds for the cone (I 4.10).

Moreover, on every open @@M10@@ of the covering we can write the commutative diagram:
@@M11@@
which induces an isomorphism between the cones; the lower homomorphism is an isomorphism in degree @@M12@@; therefore @@M13@@ is isomorphic over @@M14@@ to a bounded complex with terms locally free of finite type over @@M15@@, whose components @@M16@@ are zero for @@M17@@. Consequently, @@M18@@ has the same properties as @@M19@@, with @@M20@@ replaced by @@M21@@. We can therefore iterate the process, and it suffices to see that, after finitely many operations, the resulting complex @@M22@@ is acyclic. Here again the question is local, so it suffices to check this for @@M23@@. For each term @@M24@@, the morphisms considered above are those defining the filtration of Proposition 1.7; this filtration is finite, so, starting from @@M25@@, we reach the zero complex after finitely many operations, that is, the corresponding complex @@M26@@ is acyclic.

